\documentclass[17pt]{article}

\usepackage{cmap}	% поиск по документу
\usepackage{hyperref}	% поддержка гиперссылок
\usepackage{geometry}	% настройка геометрии документа
\usepackage{multicol}   % настройка шаблона расположения текста
\usepackage{extsizes}	% расширенная настройка размеров шрифтов
\usepackage{indentfirst}    % отступ для первого абзаца
\usepackage{tabularx}	% расширенные таблицы
\usepackage{makecell}   % настройка клеток таблиц
\usepackage{caption}    % настройка подписей к вставкам
\usepackage[table]{xcolor}  % покраска таблиц
\usepackage{graphicx}	% вставка изображений
\usepackage{wrapfig, booktabs}	% оформление таблиц
\usepackage{graphicx, caption, subcaption}
\usepackage{minted}    % вставка кода

\usepackage{color}	% настройка цвета текста - \textcolor
\usepackage{soul}	% выделение текста - \hl
\usepackage[T2A]{fontenc}	% поддержка кириллических символов
\usepackage[utf8]{inputenc} % поддержка кириллических символов
\usepackage[english,russian]{babel}	% поддержка кириллических символов
\usepackage{csquotes, biblatex}   % вывод библиографии

\usepackage{mathcmd}	% математические команды
\usepackage{amsfonts}	% математические символы 1
\usepackage{amssymb}	% математические символы 2
\usepackage{amsmath}	% математические символы 3
\usepackage{amsthm}	% работа с теоремами

% Настройка окружения теорем 1 (RU)
\newtheorem{theorem}{Теорема}[section]
\newtheorem{corollary}{Следствие}[theorem]
\newtheorem{lemma}[theorem]{Лемма}
\newtheorem{example}{Пример}[section]
\newtheorem*{remark}{Утверждение}
\addto\captionsrussian{\renewcommand\proofname{Док-во:}}

% Настройка окружения теорем 2 (EN)
%\newtheorem{theorem}{TH}[section]
%\newtheorem{corollary}{CO}[theorem]
%\newtheorem{lemma}[theorem]{LM}
%\newtheorem{example}{EX}[section]
%\addto\captionsrussian{\renewcommand\proofname{PF:}}

\renewcommand{\labelitemii}{$\circ$}
\renewcommand{\labelitemiii}{-}

\addbibresource{biblio.bib}
\graphicspath{ {./images/part1/} }

% Настройка оформления вставок кода
\setminted{xleftmargin=\parindent, fontsize=\footnotesize, linenos, tabsize=2, breaklines}
\usemintedstyle{emacs} 

% Настройка отступов
\geometry{
    a4paper,
    % total={200mm,287mm},
    left=10mm,
    right=10mm,
    top=10mm,
    bottom=20mm
}

% Заголовок
\title{\large{Методы моделирования и анализа стохастических систем\\ Лабораторная работа № 1, Часть 1}}
\author{\large{Костарев Иван, МЕНМ-150201}}
\date{весенний семестр 2025/2026}


\begin{document}

\maketitle
\tableofcontents
\newpage

\section{Инициализация ГПСЧ для X $\sim$ U[0, 1]}

Инициализируем генератор \texttt{rng}, используя конструктор \texttt{default\_rng} из библиотеки \texttt{numpy}:
\begin{minted}{python}
rng = np.random.default_rng()
\end{minted}

\section{Нахождение зависимости между мат ожиданием, дисперсией и размером выборки}

Для исследования сходимости выборочных характеристик генерируются массивы различной длины и вычисляются их математические ожидания и дисперсии.

Построим график (Рис. \ref{fig:mean_vs_N}) зависимости математического ожидания \texttt{mean} от размера выборки $N$.
\begin{figure}[htbp]
    \centering
    \includegraphics[width=0.6\textwidth]{mean_vs_N.png}
    \caption{Зависимость выборочного математического ожидания от размера выборки $N$.}
    \label{fig:mean_vs_N}
\end{figure}

Построим график (Рис. \ref{fig:var_vs_N}) зависимости дисперсии \texttt{var} от размера выборки $N$.
\begin{figure}[htbp]
    \centering
    \includegraphics[width=0.6\textwidth]{var_vs_N.png}
    \caption{Зависимость выборочной дисперсии от размера выборки $N$.}
    \label{fig:var_vs_N}
\end{figure}

\section{Построение гистограммы и вычисление погрешности (расхождение сгенерированных чисел с теоретическим значением)}

Генерируется выборка большого объема ($N = 10^8$) и строится ее нормированная гистограмма:
\begin{minted}{python}
sample_size = 10 ** 8
sample = rng.uniform(0., 1., size=sample_size)
num_bins = 32
sample_hist, bin_edges = np.histogram(sample, bins=num_bins, density=True)
\end{minted}

Построим приближенную гистограмму (Рис. \ref{fig:histogram}) равномерного распределения на основе сгенерированных данных (\texttt{sample}).
\begin{figure}[htbp]
    \centering
    \includegraphics[width=0.6\textwidth]{histogram.png}
    \caption{Приближенная гистограмма равномерного распределения $U[0, 1]$.}
    \label{fig:histogram}
\end{figure}

Определим функцию для вычисления погрешности приближения гистограммы к теоретической плотности ($p_{\text{true}} = 1$):
\begin{minted}{python}
def distribution_error(p: np.ndarray, p_true: float = 1.):
    return np.sum((p - p_true) ** 2)
\end{minted}

Построим график (Рис. \ref{fig:error_vs_N}) зависимости погрешности \texttt{distribution\_error} от размера выборки $N$.
\begin{figure}[htbp]
    \centering
    \includegraphics[width=0.6\textwidth]{error_vs_N.png}
    \caption{Зависимость погрешности распределения от размера выборки $N$.}
    \label{fig:error_vs_N}
\end{figure}

\section{Вычисление коэффициента корреляции между двумя выборками из одного распределения, но построенными разными способами}

Для оценки линейной зависимости реализуем вычисление коэффициента корреляции Пирсона:
\begin{minted}{python}
def pearson_r(x: np.ndarray, y: np.ndarray):
    return np.sum((x - x.mean()) * (y - y.mean())) / np.sqrt(np.sum((x - x.mean()) ** 2) * np.sum((y - y.mean()) ** 2))
\end{minted}

\vspace{0.5em}
\textbf{(а)} генерация двух независимых выборок.

Построим график (Рис. \ref{fig:corr_a}) зависимости коэффициента корреляции \texttt{pearson\_r} от размера выборки $N$ (вариант а).
\begin{figure}[htbp]
    \centering
    \includegraphics[width=0.6\textwidth]{corr_a.png}
    \caption{Зависимость коэффициента корреляции от размера выборки $N$ (вариант а).}
    \label{fig:corr_a}
\end{figure}

\vspace{0.5em}
\textbf{(б)} генерация единой последовательности и разделение на четные/нечетные элементы.

Построим график (Рис. \ref{fig:corr_b}) зависимости коэффициента корреляции \texttt{pearson\_r} от размера выборки $N$ (вариант б).
\begin{figure}[htbp]
    \centering
    \includegraphics[width=0.6\textwidth]{corr_b.png}
    \caption{Зависимость коэффициента корреляции от размера выборки $N$ (вариант б).}
    \label{fig:corr_b}
\end{figure}

\end{document}
